%%
% BIThesis 本科毕业设计论文模板（全英文） —— 使用 XeLaTeX 编译 The BIThesis Template for Undergraduate Thesis
% This file has no copyright assigned and is placed in the Public Domain.
%%

% 摘要若要按经管学院的要求，先英文再中文，请调换以下 abstract、abstractEn 的顺序。

\begin{abstract}
  互联网交换点（IXP）带宽转售商面临严峻挑战：客户需求随时间波动（白天1Gbps、夜间200Mbps），但传统合同分配固定容量，导致大量浪费或超额费用。本文设计并实现了一个分布式微服务控制平面，通过周期性密封拍卖机制动态分配IXP带宽，实现客户间的公平资源竞争和高效利用。系统由三个核心微服务组成：API网关接受客户的带宽竞拍（指定路径与价格），拍卖运行器在固定间隔执行密封拍卖算法计算公平分配和清算价格，遥测处理器收集数据平面流量统计并提供反馈驱动客户下一轮竞拍决策。这些服务通过Atomix分布式一致性存储（Raft协议）维护强一致的共享状态，通过Apache Kafka异步消息队列实现服务解耦与独立扩展。

  微服务架构通过解耦实现扩展性，但引入了可观测性难题：“软故障”——控制平面逻辑破损（如共识层降级、消息处理延迟）但基础设施监控仍显示正常（容器运行、CPU正常）。本文通过OpenTelemetry框架的综合可观测性设计解决此问题，对控制平面进行系统化仪器化以快速检测故障。框架采用指标（趋势检测）、分布式追踪（因果路径重建）、结构化日志（事件级证据）三支柱架构。关键技术是追踪上下文传播方案：通过在Atomix共享状态中嵌入上下文，保存了跨异步边界的W3C TraceContext，实现了从竞拍提交、拍卖清算到带宽实施的端到端请求追踪。

  实验评估通过受控故障场景（即网络流量突增期间的共识存储延迟降级）证明了框架的有效性。结果显示，可观测性框架能够实现快速的异常检测（警报在降级后数秒内触发），并能通过属性关联精准定位数据库瓶颈，使运维人员能够迅速区分网络拥塞与控制平面内部的系统故障。本系统为构建可观测的分布式微服务控制平面确立了工程实践，为设计任何需要强一致性和多服务协调的SDN应用提供了架构参考价值。
\end{abstract}

\begin{abstractEn}
  Internet Exchange Point (IXP) bandwidth resellers face significant challenges due to fluctuating customer demand. Traditional contracts allocate fixed capacity, leading to massive waste during off-peak hours and potential packet loss during surges. This thesis designs and implements a distributed microservice control plane that dynamically allocates IXP bandwidth through a periodic sealed-bid auction mechanism, enabling fair resource competition and high utilization. The system consists of three core microservices: an API Gateway ingests customer bandwidth bids, an Auction Runner executes a uniform-price clearing algorithm at scheduled intervals, and a Telemetry Processor aggregates data-plane flow statistics to drive subsequent bidding decisions. These services maintain strongly consistent shared state via the Atomix distributed consensus store (Raft protocol) and achieve asynchronous decoupling via the Apache Kafka message broker.

  While microservice decoupling provides scalability, it introduces observability challenges: ``soft failures'' occur when the control plane is logically broken (consensus degradation, message processing lag) yet infrastructure monitoring shows normalcy (containers healthy, CPU normal). This thesis addresses this gap through comprehensive observability design using the OpenTelemetry framework, systematically instrumenting the control plane to rapidly detect failures. The framework employs three complementary signals: metrics for trend detection, distributed traces for causal path reconstruction, and structured
\end{abstractEn}